\chapter{Wallpaper Groups: A User's Guide}\label{chap:physics}
%\chapter{Physics}
Corepresentations are particularly important for the crystallographic groups, also known as the wallpaper groups.
%In this chapter, we consider a slightly more general class of groups that contains them.
In this chapter we present a guide into their linear and antilinear representations
in a style that makes it easier to read for the users of these groups: crystallographers, chemists, mathematicians, physicists.
We won't prove theorems in this chapter. Instead, we state facts, enumerated using uppercase Roman numerals and
present {\em justifications} why some of these facts hold.

\section{Crystallographic Groups} \label{section:algebra_cryst}
%\textcolor{orange}{I think we need to review this section, I think it should be $\Aff_n \coloneqq \GL_n(\bR) \rtimes \bR^n$, and $E(n) \coloneqq \HO_n (\bR)\ltimes \bR^n$.}
Consider the groups of affine transformations and euclidean isometries
\[ \Aff_n \coloneqq \GL_n (\bR)\ltimes \bR^n \, , \quad \Euc_n \coloneqq \HO_n (\bR)\ltimes \bR^n \]
of the euclidean space $(\bR^n,\phi)$. Note that the subgroups $\bR^n$ are canonical and normal: they consist of all the translations.
The subgroups $\GL_n (\bR)$  and $\HO_n (\bR)$ depend on the choice of an origin:
they consist of all the transformations preserving the origin.

A subgroup $G \leq \Euc_n$ is {\em crystallographic}, if it is {\em uniform} (the quotient $\bR^n/G$ is compact)
and {\em discrete} (as a subset of $\Euc_n$). When $n=2$, we refer to $G$ as {\em a wallpaper group}.


The fundamental facts about $G$ are the three {\em Bieberbach Theorems}:
\begin{enumerate}[label=\textnormal{(\Roman*)}]
\item {\em The translation group} $T\coloneqq G \cap \bR^n$ is a lattice, i.e., $T=\bZ^n = \sum_{i=1}^n \bZ e_i$, spanned over $\bZ$ by a basis $e_1, \ldots, e_n$ of $\bR^n$.
{\em The point group} $Q\coloneqq G/T$ is finite.  
\item If $G,G'$ are two crystallographic group of the same dimension $n$ and $\psi: G \to G'$ is an isomorphism of groups,
  then $\psi$ is a conjugation in $\Aff_n$, i.e., there exists $x\in\Aff_n$ such that $\psi (g) = xgx^{-1}$ for all $g\in G$.
\item Up to conjugation in $\Aff_n$, there exist only finitely many crystallographic groups of each dimension $n$. 
\end{enumerate}  

To work with the algebraic structure of $G$,
we choose a representative $\breve{q}\in G$
for each element $q\in Q$, 
The finite set of all these representatives $\breve{Q} \coloneqq \{ \breve{q}\mid q\in Q\}$ is called {\em the transversal}.
The point group $Q$ naturally acts on the lattice $T$ via conjugation: $q(t) \coloneqq \breve{q}t\breve{q}^{-1}$.

\begin{just}
%  At this stage, one needs to verify two things: independence of the choice of $\breve{q}$ and
  The action axiom is $q_1(q_2 (t)) = q_1q_2 (t)$. Its verification goes through because $T$ is abelian.
  Naturally in this context means that the action does not depend on the choice of origin and the choice of transversal.
  Both are true because the action can be written
as $q(t)(x) = q(t(x)) - q(x)$ for any $x\in \bR^n$.
\end{just}  

%Thinking of $T=\bZ^n$ inside $\bR^n$, allows us to pinpoint the group $Q$ as a subgroup of $\GL_n (\bZ)\cap \HO_n (\bR)$.
It is convenient to write all the matrices in the basis $e_1, \ldots, e_n$.
Let $\Phi = \begin{pmatrix}\phi (e_i,e_j)\end{pmatrix}$ be the Gram  matrix of this basis.
Then $Q$ is a subgroup of $\GL_n (\bZ)\cap \HO_n (\bR, \Phi) = \{ x\in \GL_n (\bZ) \mid x^{T}\Phi x=\Phi\}$.
The transversal element of $q\in Q$ can be written as $\breve{q} = [q|a]$, which is
a transformation $\breve{q} (x) = q(x) + a$. 


Each point $x\in \bR^n$ gives also {\em the isotropy group} $G_x \coloneqq \{ g\in G \mid g(x) = x\}$, often called {\em the site-symmetry group}.
Each $G_x$ is naturally a subgroup of $Q$ under the map $G_x \to G \twoheadrightarrow Q$. Every isotropy group gives a semidirect product subgroup
$G_x \ltimes T$ inside $G$. The next statement is a standard algebra fact.
\begin{enumerate}[resume, label=\textnormal{(\Roman*)}]
\item The following five statements about $G$ are equivalent:
  \begin{itemize}
\item   $G$ is a semidirect product $Q\ltimes T$.
\item  The following exact sequence splits:
  \begin{equation} \label{sequence:TGQ}
    1\to T\to G \xrightarrow{\varpi} Q \to 1 \, .
    \end{equation}
\item  There exists a point $x\in \bR^n$ with $G_x=Q$.
\item One can choose a transversal $\breve{Q}$, which is a subgroup. 
\item With a suitable choice of the origin, one can choose a transversal $\breve{Q}$, where $\breve{q} = [q|0]$ for each $q\in Q$. 
\end{itemize}
  \end{enumerate}

Let us choose an origin from now on.
Then every element $g\in G$ can be written as $g =[q|a]$ for unique $q\in Q$ and $a\in \bR^n$. It is worth writing the formulas for multiplication and inverse in $G$ and the conjugation action on $T$ in these terms:
\begin{equation} \label{mult_wallpaper}
  [q_1|a_1][q_2|a_2] =[q_1q_2|q_1(a_2)+a_1]\, , \quad 
  [q|a]^{-1}= [q^{-1}|-q^{-1}(a)]\, , \quad  [q|a]t[q|a]^{-1} = q(t) \, .
\end{equation}  
\begin{just}
  The first formula is verified by applying both sides to $x\in\bR^n$.
The second and third formulas are deduced from the first one.
\end{just}  

We need a description of the centre $Z=Z(G) \coloneqq \{ g\in G \mid \, \forall h\in G \; gh=hg\}$.
Let us restrict the exact sequence~\eqref{sequence:TGQ} to the centre:
\begin{equation} \label{sequence:TGQ_centre}
  1 \to Z \cap T \to Z \xrightarrow{\varpi} Z(Q) \, .
\end{equation}  
%The following two facts follow from~\eqref{mult_wallpaper}:
\begin{enumerate}[resume, label=\textnormal{(\Roman*)}]
\item In the exact sequence~\eqref{sequence:TGQ_centre}, the kernel $Z\cap T$ is equal to $T^Q = \{ t\in T \mid \, \forall q\in Q \; q \cdot t = t \}$
  and the image of $\varpi$ is trivial.  Therefore, $Z(G) =T^Q$.
%\item $q\in Z(Q)$ belongs to the image of $\varpi$ if and only if $q$ admits a lifting $[q|x]\in G$ with $x\in (\bR^n)^Q$.
\end{enumerate}
\begin{just} The exact sequence is the restriction of the map $\varpi: G \to Q$ to $Z$. Write $t\in T$ as $[I|t]\in T$ and use~\eqref{mult_wallpaper}:
  \[ [q|x][I|t] = [I|t][q|x] \ \Leftrightarrow \  [q| q(t)+x]= [q| t+x] \ \Leftrightarrow \ q(t) = t \, .\]
  This proves both facts. For $t\in T$ to belong to $Z$, we need $q(t)=t$ for all $q\in Q$. And for
$[q|x]\in Z$ we need $q(t)=t$ for all $t\in T$, forcing $q=I$.  
%Ditto for the second fact. Fix $q\in Z(Q)$. For all $[p|y]\in G$,
%  \[ [q|x][p|y] = [p|y][q|x] \ \Leftrightarrow \  [qp| q(y)+x]= [pq| y+p(x)] \ \leftrightarrow \ q(t) = t \, .\]
\end{just}

Another algebraic piece of the puzzle is {\em the original cocycle}, 
    \[ \beta_{o} \in \HZ^1 (Q,\bR^n/T), \quad  [ q | \beta_o (q) ] = \varpi^{-1} (q) = G \cap [ q | \bR^n ]\, . \]
    To understand it better, choose an origin $o$ and write the transversal elements explicitly as $\breve{q} = [q| a_q]$.
    Then $\beta_o (q) = a_q+T$, so that $\beta_o$ essentially describes the fibres of the map $\varpi : G \to Q$.
    It depends on the origin but it becomes canonical, once an origin is chosen. Hence the name.

    Notice that the original cocycle maps to
    the 2-cohomology responsible for the original extension~\eqref{sequence:TGQ} under the maps
    \[ \HZ^1 (Q,\bR^n/T) \twoheadrightarrow \HH^1 (Q,\bR^n/T) \xrightarrow{\cong} \HH^2 (Q,T) \, . \]

    Finally consider {\em the original lattice} $T_o$, generated by $T$ and all $a_q$, $q\in Q$.
    It is always a free abelian group but it may or may not be a discrete subgroup $\bR^n$. It depends on the origin but does not depend on the choice of transversal.
  \begin{enumerate}[resume, label=\textnormal{(\Roman*)}]
  \item The original lattice is a $Q$-submodule of $\bR^n$ containing $T$. One can choose an origin so that
    $T_o \subseteq \frac{1}{|Q|} T$.
  \end{enumerate}
  \begin{just}
    By definition, $T_o$ is an abelian group and $T\subseteq T_o \subset \bR^n$. Pick $s,t\in Q$. By~\eqref{mult_wallpaper},
    $q(a_s)+a_q \in T_o$ for all $q,s\in Q$. Thus, $q(a_s) \in T_o$.

    The second statement follows immediately from the transfer principle in group cohomology but we give an elementary recipe.
    Start with any origin $o$, then move it to $o+ \frac{1}{|Q|} \sum_{q\in Q} a_q$. The reader can easily check that the new origin satisfies
    the requirements.
\end{just}

Notice that $T_o$ is not a $\GL_n(\bZ)$-module, in general. But $\widetilde{T}\coloneqq \frac1{|Q|}T$ is a $\GL_n(\bZ)$-module.
  Let us call an origin {\em rational} if $T_o \subseteq \widetilde{T}$. With a rational origin, we can always realise the group $G$ as a finite index subgroup in a split group
  \[ G \leq Q \ltimes T_o \leq Q \ltimes \widetilde{T} \, . \]

  
\section{Automorphism Groups}

Let us fix a rational origin $o$ throughout this section.

The second Bieberbach Theorem yields a description of the automorphism groups $\Aut (G)$ and $\Out (G)$ via
the normaliser $N_{\Aff_n} (G) = \{ x \in \Aff_n \mid \forall g\in G \ xgx^{-1}\in G \}$ and
the centraliser $C_{\Aff_n} (G) = \{ x \in \Aff_n \mid \forall g\in G \ xgx^{-1}=g \}$:
\begin{equation} \label{formula:AUT1}
  \Aut (G)  = N_{\Aff_n} (G) /C_{\Aff_n} (G), \quad
  \Out (G)  = N_{\Aff_n} (G) /G \cdot C_{\Aff_n} (G).
\end{equation}  

We would like another, more efficient description.
Note that $\varphi (T) =T$ for any automorphism $\varphi \in \Aut (G)$ as follows from the second Bieberbach Theorem.
Consequently, any automorphism $\varphi \in \Aut (G)$ has three components
\begin{equation} \label{formula:AUT0}
  \varphi [q|x] = [ \varphi_Q (q) | \varphi_T (x) + \varphi_m (q,x)] \, , 
\end{equation}
where $\varphi_Q\in \Aut (Q)$ and the remaining components
$\varphi_T$ and $\varphi_m$ may be represented by an element $\Aut (T_o) = \GL_n (\bZ)$ and a map $Q \times T_o \to T_o$.
Let us make it precise. 
  \begin{enumerate}[resume, label=\textnormal{(\Roman*)}]
  \item For every $\varphi \in \Aut (G)$ there exist unique $A\in N_{\GL_n (\bZ)}(Q)$ and $\alpha \in \HZ^1 (Q,T_o)$
  such that   
\begin{equation} \label{formula:AUT}
  \varphi [q|x] = [ AqA^{-1} | A(x + \alpha (q))] \, .
\end{equation}
  \end{enumerate}
  \begin{just}
    By the second Bieberbach Theorem, $\varphi$ is a conjugation by $[A|a]\in \Aff_n$.  Using~\eqref{mult_wallpaper}, we can write
    \[
    \varphi [q|x] =
            [ A|a] [q|x][A|a]^{-1}  =
            [ Aq | Ax +a][A^{-1}| - A^{-1}a]  =
                       [ AqA^{-1} |Ax +a - AqA^{-1}a] \, .
    \]
    Since $\phi [I|t] = [I|At]$ for all $t\in T$, we conclude that $A\in \GL_n (\bZ)$.
    Since $AqA^{-1}\in Q$ for all $q\in Q$, $A$ must actually belong to the normaliser $N_{\GL_n (\bZ)}(Q)$.
    Finally, the function 
    \[ \alpha : Q \to \bR^n, \quad \alpha (q) \coloneqq A^{-1}a - qA^{-1}a\]
    takes values in $T_o$. It is a 1-cocycle, by a direct calculation, i.e.,  $d \alpha \equiv 0_T$ where
    \[d \alpha (q,s) \coloneqq q\cdot \alpha (s) - \alpha (qs)+ \alpha (q) \, . \]
    Finally, suppose $(A_1,\alpha_1)$ and $(A_2,\alpha_2)$ represent the same automorphism.
    Then
    \[
    (A,\alpha) \coloneqq (A_1,\alpha_1)^{-1}\cdot(A_2,\alpha_2) = (A_1^{-1}A_2, A_1^{-1}\alpha_2 - A_1^{-1}\alpha_1)^{-1} \in N_{\GL_n (\bZ)}(Q) \ltimes \HZ^1 (Q,T_o)
    \] 
    represents the trivial automorphism $\varphi$. Applying~\eqref{formula:AUT}, $A(x+\alpha(I)) = x+\alpha(I)$ for all $x\in T$.
    Thus, $A=I$.
    Applying~\eqref{formula:AUT} again, $a_q+\alpha(q) =a_q$ for all $q\in Q$. Hence, $\alpha =0$. 
  \end{just}
  
  So far  we have realised $\Aut (G)$ as a subgroup of the semidirect product
  $N_{\GL_n (\bZ)}(Q) \ltimes \HZ^1 (Q,T_o)$. To describe it more precisely,
  we need little more algebraic yoga. Notice that the original cocycle $\beta_o$ has values in $T_o/T$.
  We can think that $\beta_o \in \HZ^1 (Q, T_o/T)$.
  We want to push $\alpha \in \HZ^1 (Q, T_o)$ forward to the same group:
  let $\overline{\alpha} \coloneqq (T_o \to T_o/T)_\ast (\alpha) \in \HZ^1 (Q, T_o/T)$.
%    Finally, we need to pinpoint the following right $Q$-action on $\HZ^1 (Q, T_o/T)$
%    \[ \gamma^q : s \mapsto q^{-1} \cdot \gamma (qsq^{-1}) \, . \]
        Finally, we need to pinpoint the following right $N_{\GL_n (\bZ)}(Q)$-action on $\HZ^1 (Q, \widetilde{T}/T)$
    \[ \gamma^A : s \mapsto A^{-1} \cdot \gamma (AsA^{-1}) \, . \]
Under this action, $\HZ^1 (Q, T_o/T)$ is a $Q$-submodule of $\HZ^1 (Q, \widetilde{T}/T)$.
Observe that the induced $Q$-action on $\HH^1 (Q, \widetilde{T}/T)$ is trivial.
Therefore,  $N_{\GL_n (\bZ)}(Q, o)$,
the stabiliser of the original cohomology class $[\beta_o]\in \HH^1 (Q,\widetilde{T}/T)$ inside the normaliser $N_{\GL_n (\bZ)}(Q)$,
contains $Q$.
\begin{enumerate}[resume, label=\textnormal{(\Roman*)}]
\item
  A pair $(A,\alpha)\in N_{\GL_n (\bZ)}(Q)\times \HZ^1 (Q,T_o)$ defines an automorphism $\varphi = (A,\alpha)$ by the formula~\eqref{formula:AUT}
  if and only if
  \[A \in N_{\GL_n (\bZ)}(Q, o) \ \mbox{ and } \ \beta_o^A = \beta_o + \overline{\alpha} \, . \]
   Moreover,
%  the full automorphism group is a semidirect product 
%  \[ \Aut (G)=   N_{\GL_n (\bZ)}(Q,\beta) \ltimes \HZ^1 (Q,T)\, .\]
    the full automorphism group fits into a short exact sequence 
    \begin{equation} \label{Physics_AutSequence}
      1 \longrightarrow \HZ^1 (Q,T) \longrightarrow \Aut (G) \longrightarrow   N_{\GL_n (\bZ)}(Q,o) \longrightarrow 1\, .
    \end{equation} 
\end{enumerate}

\begin{just}
  Pick $(A,\alpha)\in N_{\GL_n (\bZ)}(Q)\times \HZ^1 (Q,T_o)$.
  A straightforward algebraic check tells that  the formula~\eqref{formula:AUT} defines an injective homomorphism $G\to \Aff_n$.
  Thus, we just need to control its image. Using the transversal, observe that
  $(A,\alpha)$ being an automorphism is equivalent to the condition
  \[ \forall q\in Q \quad (A,\alpha) [q| a_q] = [ AqA^{-1} | A(a_q + \alpha (q)) ] \in G \, . \]
  Clearly, this condition is equivalent to another condition with the original cocycle
  \[ \forall q\in Q \quad Aa_q +A \alpha(q)\in  a_{AqA^{-1}} + T  = \beta_o (AqA^{-1}) \, . \]
  Modulo $T$, this is equivalent to
  \[a_q + \alpha(q)+ T = A^{-1}\beta_o (AqA^{-1}) \quad \mbox{or} \quad
\beta_o + \overline{\alpha} = \beta_o^A \, ,\]
as stated. Now observe that $\overline{\alpha}$ is a 1-coboundary. Hence, 
$[\beta_o]  = [\beta_o]^A$ and $A\in N_{\GL_n (\bZ)}(Q,o)$. It follows that
\[
f: (A,\alpha) \mapsto A, \quad \Aut (G) \longrightarrow   N_{\GL_n (\bZ)}(Q,o) 
\]
is a surjective group homomorphism. Its kernel consists of
of pairs $(I,\alpha)$. But $(I,\alpha)$ defines an automorphism of $G$ if and only if $\overline{\alpha}=0$. In other words,
$\alpha \in \HZ^1 (Q,T)$. Therefore,  $\HZ^1 (Q,T)$ is the kernel of $f$.
\end{just}


As observed above, $Q$ is a subgroup of $N_{\GL_n (\bZ)}(Q, o)$.
It must be normal since $Q$ is normal in $N_{\GL_n (\bZ)}(Q)$.
The group of inner automorphisms $\Inn (G)$ is the image of $G$ inside $\Aut (G)$.
It can be described in terms of the exact sequence~\eqref{Physics_AutSequence} as
\begin{equation} \label{Physics_InnSequence}
  1 \longrightarrow \HB^1 (Q,T) \longrightarrow \Inn (G) \longrightarrow  Q \longrightarrow 1\, .
\end{equation} 
This yields a short exact sequence that describes $\Out (G)$ 
    \begin{equation} \label{Physics_OutSequence}
      1 \to \HH^1 (Q,T) \to \Out (G) \to N(Q, o) \to 1\, , \ \mbox{ where } N(Q, o) \coloneqq N_{\GL_n (\bZ)}(Q, o)/Q \, .
      \end{equation}
    If the sequence~\eqref{sequence:TGQ} splits, then three simplifications take place:
    $\beta =0$, $N_{\GL_n (\bZ)}(Q, o) = N_{\GL_n (\bZ)}(Q)$ and the sequence~\eqref{Physics_OutSequence} also splits.
%It may also split for other reasons: for instance, it always splits in dimension 2.    
We tabulate all the necessary data in the 2-dimensional case in Appendix~\ref{appendix_wallpaper}.

\section{Magnetic Wallpaper Groups} \label{sect:are_there}

A \emph{magnetic group} is subgroup $M \subset \Euc_n \times C_2$, such that the image $\pi_1(M)$ under projection onto the first factor $\Euc_n$ is crystallographic. Two magnetic groups are considered \emph{equivalent} if conjugation by an element of $\Aff_n$ %on $\Euc_n \times C_2$
takes one to the other.


Now magnetic groups come in three classes, corresponding to the three possible classes of subgroup of a direct product $\Euc_n \times C_2$:
\begin{itemize}
  \item Ordinary crystallographic groups, $M = G \times 1 < \Euc_n \times C_2$,
  \item ``Grey'' groups, $M = G \times C_2 < \Euc_n \times C_2$, where $G$ is crystallographic, and
  \item ``Black-and-white'' or ``two-colour'' groups
    %\[
    $M = \{(g, \phi(g)) \mid g \in G\}$
    %\]
  for some surjective homomorphism $\phi \colon G \twoheadrightarrow C_2$.
\end{itemize} 


Observe that the data which determines a black-and-white group is a pair $(G, \phi \colon G \twoheadrightarrow C_2)$ where $G$ is crystallographic, which is the same as the data of the pair $(G,\ker(\phi))$, i.e. a crystallographic group and an index two subgroup. Why are they called black-and-white? If $G$ is $2$-dimensional, a wallpaper group, then $G$ is the symmetry group of some wallpaper pattern. The map $\phi$ corresponds to a colouring of this pattern in black and white so that $g$ preserves the colouring if $\phi(g) = 1$ and swaps the colours if $\phi(g) \neq 1$, see \cite{Schat}.

We can also give a more algebraic point of view, which makes them easier to classify. Let us define a \emph{Magnetic group} (with a capital $M$) to be a $C_2$-graded group (meaning a group $G$ equipped with a homomorphism $\varphi \colon G \rightarrow C_2$) such that the ungraded subgroup $H = \ker(\varphi)$ is crystallographic. Two Magnetic groups are the same (are \emph{equivalent}) if they are isomorphic as $C_2$-graded groups: there is some group isomorphism $\psi \colon G_1 \rightarrow G_2$ such that $\psi(H_1) = H_2$.

From a magnetic grey or black-and-white group $M \subset \Euc_n \times C_2$, one obtains a Magnetic group $(M, \pi_2 \colon M \rightarrow C_2)$ via projection to the second factor $C_2$. The key property is that every Magnetic group is magnetic.

\begin{enumerate}[label=\textnormal{(\Roman*)}]
\item
  Suppose $H$ is a crystallographic group and $H \leq G$ is a Real structure. Then $G$ is either grey, i.e.,
  $G \cong H \times C_2$, or black-and-white, i.e., $G$ is also crystallographic (in the way compatible to the crystallographic structure of $H$) of the same dimension as $H$.
\end{enumerate}

Note that $G$ cannot be both grey and black-and-white: $H \times C_2$, for $H$ crystallographic, has torsion elements in $Z(H \times C_2)$, which is impossible for crystallographic groups, cf. Fact~\ref{section:algebra_cryst}(V).

\begin{just}
  Let $T\subset H$ be the translation lattice and $Q$ the point group of $H$.
Set $V\coloneqq (\bR^n)^Q$. Observe first that
\[
C_{\Aff_n}(H) \;=\; C_{\Euc_n}(H) \;=\; \{\,[I|v] \mid v\in V\,\}.
\]
Indeed, if $[A|b]$ commutes with every $[I|t]$, $t\in T$, then $At=t$ for all $t\in T$; since $T$ spans $\bR^n$, $A=I$.
Then $[I|b]$ commutes with $[B|c]\in H$ if and only if $Bb=b$. In particular, $C_{\Aff_n}(H)$ is an $\bR$-vector space, so for any $C_2$-action $\HH^1 (C_2, C_{\Aff_n}(H))=0$ -- this is the heart of the proof below.

Pick $g\in G\setminus H$ and consider the automorphism $\gamma_g\in\Aut(H)$, $\gamma_g(h)=ghg^{-1}$.
By the second Bieberbach theorem, $\gamma_g=\gamma_x$ for some $x\in\Aff_n$.
This $x$ is unique up to right multiplication by $C_{\Aff_n}(H)$: one may replace it by $x[I|v]$ for any $v\in V$.
Let $A$ be the linear part of $x$. Since $x$ normalises $H$, conjugation by $x$ preserves $C_{\Aff_n}(H)$, so $AV=V$, and
\[
[I|v]\,x \;=\; x\,[I|A^{-1}v] \qquad (v\in V).
\]
We hope to choose $v$ so that
\[
\psi : G\to\Aff_n,\qquad \psi(h)=h,\quad \psi(hg)=h\,x[I|v]\qquad (h\in H)
\]
is a homomorphism. Note that $G$ is generated by $H$ and $g$ subject to the following relations:
\[ h_1h_2 = h_1h_2 \in H, \ \ ghg^{-1}=\gamma_g(h) \ \mbox{ and } \ g^2=g^2\in H \, . \]
Since $\gamma_{x[I|v]}=\gamma_x=\gamma_g$ on $H$,
$\psi$ is a well-defined homomorphism if and only if
\[
g^2 \;=\; (x[I|v])^2 \;=\; x^2\,[I|A^{-1}v+v]\,.
\]
Now $g^2$ and $x^2$ induce the same automorphism of $H$, so $x^{-2}g^2\in C_{\Aff_n}(H)$, i.e.\ $g^2=x^2[I|w]$ for some $w\in V$.
Since $g^2\in H$ commutes with $g$ and $\gamma_x=\gamma_g$, we have $xg^2=\gamma_x(g^2)\,x=\gamma_g(g^2)\,x=g^2x$, hence
\[
g^2[I|-w]\,x \;=\; x^2\cdot x \;=\; x\cdot x^2 \;=\; x\,g^2[I|-w] \;=\; g^2\,x\,[I|-w]\,,
\]
so $[I|-w]x=x[I|-w]$, i.e.\ $Aw=w$.
Let $v\coloneqq\tfrac12 w$. Then $A^{-1}v=v$ and
\[
(x[I|v])^2 \;=\; x^2\,[I|A^{-1}v+v] \;=\; x^2\,[I|w] \;=\; g^2 ,
\]
so $\psi$ is a well-defined homomorphism, injective on $H$.

Let $K\coloneqq\ker\psi$. Since $K\cap H=1$ and $[G:H]=2$, either $K=1$ or $K=\langle k\rangle$ has order $2$.

If $|K|=2$, then $k\notin H$ and, being a normal subgroup of order $2$, $\langle k\rangle$ is central; hence
$G=H\times\langle k\rangle\cong H\times C_2$ with the $C_2$-factor in $G\setminus H$, i.e.\ $G$ is grey.

If $K=1$, then $\Gamma\coloneqq\psi(G)\subset\Aff_n$ contains $T$ as a normal subgroup of finite index.
Its group $Q_\Gamma\subset\GL_n(\bR)$ of linear parts is finite, since $T$ lies in the kernel of the linear-part map.
Averaging, $\langle u,v\rangle'\coloneqq\sum_{B\in Q_\Gamma}\langle Bu,Bv\rangle$ is a $Q_\Gamma$-invariant inner product, and so there is some $S\in\GL_n(\bR)$ with $\langle u,v\rangle'=\langle Su,Sv\rangle$. Then $SQ_\Gamma S^{-1}\subset\mathrm{O}(n)$, and since $[S|0][B|b][S|0]^{-1}=[SBS^{-1}|Sb]$, the group $\Gamma' \coloneqq [S|0]\,\Gamma\,[S|0]^{-1}$ lies in $\Euc_n$. As the crystallographic subgroup $H' = [S|0]\,\Gamma\,[S|0]^{-1}$ has index two in $\Gamma'$, we see that $G$ is crystallographic of the same dimension as $H$.
\end{just} 

The following is easy to see using (I) above and Bieberbach (II).
%This essentially tells us that the magnetic equivalence of magnetic groups is precisely the equivalence of Real structures.

\begin{enumerate}[resume, label=\textnormal{(\Roman*)}]
\item
  Up to equivalence, every Magnetic group $(G, \varphi \colon G \rightarrow C_2)$ is of the form $(M, \pi_2)$ for some magnetic group $M$.
  If two magnetic groups $M, M' \subset \Euc_n \times C_2$ are equivalent, they are of the same class.
  Two ordinary magnetic groups $M, M' \subset \Euc_n \times 1$ are equivalent
  if and only if they are isomorphic.
Two grey or black-and-white magnetic groups $M, M' \subset \Euc_n \times C_2$ are equivalent
if and only if there exists an isomorphism $\psi: M \to M'$ that agrees with the second projection, that is,
$\pi_2' \psi = \pi_2$.
\end{enumerate}

The main consequence is our notation of a Magnetic group contains is essentially the same as that of a magnetic group, and this makes the classification of magnetic groups more of an algebraic problem, for example as one can now use the theory from Chapter \ref{chapter:realstructures}.

We are interested in magnetic groups in dimension $2$ -- how many are there? It was shown by Heesch \cite{Heesch} in 1929 that there are 80. See also \cite{Schat} for a mathematical treatment and \cite{PaZa, Souv, Zamo} for generalisations. Similar counts are known in dimensions 3 and 4 \cite{Souv} but we won't address them here.

In dimension 2 the count goes like this:
  \begin{itemize}
  \item 17 ordinary wallpaper groups $G$,
  \item 17 grey groups $G\times C_2$ for each wallpaper group $G$,
  \item 46 black-and-white groups $H\leq G$, where $G$ is a wallpaper group, $H$ an index two subgroup.
  \end{itemize}

We are interested in the corepresentations of magnetic groups, and as we need to work with these groups explicitly, we review this classification in Appendix \ref{appendix_wallpaper}. In fact, we derive this classification in two ways: by counting the index two subgroups, and by the more abstract classification of Real structures on wallpaper groups, using the results of Chapter \ref{chapter:realstructures} and the Magnetic groups point of view.

The interesting case are the black-and-white groups. We tabulate these in Table \ref{table:wallpaper3}.

\section{Classification of Irreducible Corepresentations}\label{sect:cliffordtheory}
Consider a crystallographic group $H$ and a Real structure $H\leq G$.
In this section we give a classification of irreducible corepresentations of $G$ via
a version of Clifford Theory. 

Let us do a quick recap of its classical version, where $K$ is a group and $N$ is its normal subgroup.
There is a left action of $K$ and $K/N$ on the set ${\Irr_{\bC}(N)}$ by conjugation:
\begin{equation} \label{Cliff_1a}
  {}^{k}(V,\rho) \coloneqq (V, n \mapsto \rho (k^{-1}nk)) \ \mbox{ where } \ n\in N, \, k\in K.
  \end{equation} 

Clifford Theory can be viewed as a correspondence between certain subsets of ${\Irr_{\bC}(K)}$ and the $K$-orbits of ${\Irr_{\bC}(N)}$ along the following lines.

  \begin{enumerate}[label=\textnormal{(\Roman*)}]
  \item Given $V\in {\Irr_{\bC}(K)}$, its restriction to $N$ is a sum of irreducible representations over the orbit of a single $W\in {\Irr_{\bC}(N)}$:
\begin{equation} \label{Cliff_1b}
  V|_N \cong \left(\oplus_{U \in {}^{K}W} U\right)^{\oplus m}\, . 
\end{equation} 
\item Given $(W,\rho) \in {\Irr_{\bC}(N)}$, consider the stabiliser of its isomorphism class $K_W$. Note that $K_W$ contains $N$:
$\rho(n)$ is an isomorphism ${}^n W \to W$. Now for each $k\in K_W \setminus N$, pick an isomorphism of $N$-modules $\rho(k) : {}^k W \to W$.
The newly extended map $\rho : K_W \to \GL (W)$ is a projective representation, in general, defining
%a cochain and the cocycle:
a cocycle and the cohomology class:
\begin{equation} \label{Cliff_1c}
d\rho \in \HZ^2 (K_W , \bC^\times), \ d \rho (t,s) = \rho (t) \rho (s) \rho (ts)^{-1}, \quad
c_W = [ d \rho]^{-1} \in \HH^2 (K_W , \bC^\times) \, .
\end{equation}  
\item The cohomology class $c_W$ does not depend on the choice of $\rho$, only on $W$.
It represents an obstruction to turning $W$ into a representation of $K_W$.  
In particular, its restriction to $N$ is trivial.
This may or may not lift\footnote{There is no naive exact sequence connecting cohomology of $N$, $K_W$ and $K_W/N$. The relations is governed by the Lyndon-Hochschild-Serre spectral sequence. This yields the seven-term inflation-restriction sequence, where one can see the obstructions to existence and uniqueness of $\overline{c_W}$.}
to a class $\overline{c_W} \in  \HH^2 (K_W/N , \bC^\times)$.
\item Suppose such a lift $\overline{c_W}$ exists (which in general will not be unique).
  The key set of representations is ${\Irr_{W}(K_W/N)}$, the set of irreducible finite-dimensional $\overline{c_W}$-projective representations of $K_W/N$. For any $X\in {\Irr_{W}(K_W/N)}$, the tensor product $W\otimes X$
  is a linear representation of $K_W$ -- the cocycles of the two tensor multiplicands cancel each other out.
\item Finally given $W$ and $\overline{c_W}$ as above, there is a one-to-one correspondence between
  ${\Irr_{W}(K_W/N)}$ and those $V\in {\Irr_{\bC}(K)}$ that contain $W$ in their restriction to $N$, defined by
\begin{equation} \label{Cliff_1d}
  X \mapsto \Ind_{K_W}^K W \otimes X, \, \qquad V \mapsto \Hom_N(W,V), \,
\end{equation}  
where $\Hom_N(W,V)$ is a $c_W$-projective representation of $K_W/N$ via $k * f \coloneqq \rho_V(k) \circ f \circ \rho(k)^{-1}$.
\end{enumerate}

Let us now readjust everything to Real structures and corepresentations. We have a crystallographic group $H$ sitting inside a Real structure $G$ and $T$ is the lattice of $H$. The pair $(K,N)$ above will be replaced with the pair $(G,T)$. Note that the translation subgroup $T$ remains normal in $G$ because $G$ acts on $H$ by automorphisms and each automorphism preserves translation by the second Bieberbach Theorem.

We fix a transversal $\breve{S}$ of the finite group $S \coloneqq G/T$, a section of the map $G \twoheadrightarrow S$, such that $\breve{1_S} = 1_G$. Because $T$ is ungraded, the group $S$ inherits the grading $\pi : T \to C_2$ from $G$.
Since $T$ is abelian, all representations ${\Irr_{\bC}(T)}$ are one dimensional, i.e.,
they are characters, elements of $T^\vee \coloneqq \hom (T,\bC^\times)$.
The group $S$ acts on $T$ by conjugation $t^s = \breve{s}^{-1} t \breve{s}$,
which translates into {\em the left Real action} on $T^\vee$: for $s\in S$ and $\chi\in T^\vee$,
\begin{equation} \label{Cliff_2a}
{}^{s}\chi (t)  \coloneqq {}^{\pi (s)} \chi (\breve{s}^{-1}t\breve{s}) = \begin{cases}
    \overline{\chi (\breve{s}^{-1}t\breve{s})}  , & \text{if $s$ odd,}\\
    \chi (\breve{s}^{-1}t\breve{s}) , & \text{if $s$ even.}\\
    \end{cases}
  \end{equation}  
This is compatible with the natural action of $S$ on $\Irr_{\bC}(T)$ as considered in Chapter \ref{chapter:twelvefoldway}, i.e.
\[
  s * \bC_{\chi} = \bC_{{}^{s} \chi},
\]
where $\bC_{\chi} \in \Irr_{\bC}(T)$ is $\bC$ with action of $T$ by multiplication by $\chi$.

Let us now describe Clifford Theory in this context, which allows one to describe ${\Irr_{\bC}^{\hspace{0.1em} \star}(G)}$ in terms of the $S$-orbits on $T^\vee$ and projective antilinear representations of quotients of $S$.

  \begin{enumerate}[resume, label=\textnormal{(\Roman*)}]
  \item Given $V\in {\Irr_{\bC}^{\hspace{0.1em} \star}(G)}$, its restriction to $T$ is a sum of characters over a single orbit, i.e.,
\begin{equation} \label{Cliff_2b}
  V|_T \cong (\oplus_{\chi \in S \cdot \mu} \bC_\chi) ^{\oplus m}\, .
\end{equation}
\end{enumerate}

\begin{just}
%  First, pick $V\in {\Irr_{\bC}^{\hspace{0.1em} \star}(H)}$.
  Elements of $T$ give a commuting set of linear operators on the finite dimensional $\bC$-vector space $V$.
  Hence, they have a common eigenvector, which spans $\bC_\chi$ for some $\chi\in T^\vee$.
  Then $V_\chi$, the sum of all $T$-submodules of $V$ isomorphic to $\bC_\chi$, is not zero. Pick $s\in S$ and observe the $T$-action on $\breve{s}V_\chi$:
\begin{equation} \label{Cliff_3a}
  t\breve{s} v = \breve{s}(\breve{s}^{-1}t\breve{s})v = \breve{s} \chi (\breve{s}^{-1}t\breve{s}) v =
  \begin{cases}
    \overline{\chi (\breve{s}^{-1}t\breve{s})}   \breve{s}v , & \text{if $s$ odd,}\\
    \chi (\breve{s}^{-1}t\breve{s})          \breve{s}v , & \text{if $s$ even.}\\
  \end{cases}
\end{equation}
This justifies~\eqref{Cliff_2a}. Furthermore, $\sum_{s\in S} \breve{s}V_\chi$ is a non-zero  $H$-submodule. Since $V$ is irreducible,
it must be $V$: this justifies~\eqref{Cliff_2b}.
\end{just}

\begin{enumerate}[resume, label=\textnormal{(\Roman*)}]
\item Given a character $\chi\in T^\vee$, we need its stabilisers $S_\chi$ and $G_\chi$ under the action~\eqref{Cliff_2a}.
Similarly to~\eqref{Cliff_1c}, we have a cochain and the cocycle, which, in our case, can be always reduced to the smaller group $S_\chi$ and found by an explicit formula. To do that, we consider the cocycle defining the product on the transversal, $\gamma \in \HZ^2 (S, T)$, which is defined by the formula
\begin{equation} \label{Cliff_2c_0}
\gamma (h,k) \cdot \breve{g}  = \breve{h} \cdot \breve{k}
\end{equation}
for $h,k\in S$, where $g \coloneqq hk$. Then $d\chi$ defined by
\begin{equation} \label{Cliff_2c}
d\chi \in \HZ^2 (S_\chi , \bC^\times), \  
d \chi (h,k) = \chi (\gamma (h,k))^{-1}, \quad
c_\chi \coloneqq [ d \chi]^{-1} \in \HH^2 (S_\chi , \bC^\times)\, ,
\end{equation}  
is well-defined, where $\bC^{\times}$ is \emph{not the trivial $S_{\chi}$-module} but rather has the action where elements of $(S_{\chi})_0$ act as the identity and all other elements of $S_{\chi}$ act by complex conjugation. Of course, if $S_{\chi}$ happens to be trivially graded, then $\bC^{\times}$ will be the trivial $S_{\chi}$-module.
\item In analogy with the classical picture, taking $W = \bC_{\chi}$, so $G_{\chi} = K_W$, the image of $c_\chi$ under the inflation map $\HH^2(S_\chi, \bC^{\times}) \rightarrow \HH^2(G_{\chi}, \bC^{\times})$ is the the class $c_W$ defined by
\begin{equation}\label{Cliff_1c_new}
d\rho \in \HZ^2 (G_{\chi} , \bC^\times), \ d \rho (t,s) \rho (ts) = \rho (t) \rho (s), \quad
c_W = [ d \rho]^{-1} \in \HH^2 (G_{\chi} , \bC^\times) \, 
\end{equation}  
where for each $g \in G_{\chi}$ we have chosen an isomorphism $\rho(g) \colon g * W \xrightarrow{\sim} W$, and similarly $\bC^{\times}$ has this potentially non-trivial action of $G_{\chi}$.
\item The class $c_\chi$ does not depend on the choice of transversal, only on the character $\chi$.
  If $G_\chi$ is a semidirect product of $S_\chi \ltimes T$,
    then the class $c_\chi \in \HH^2 (S_\chi , \bC^\times)$ is automatically trivial. 
\end{enumerate}

\begin{just}    
  To see that~\eqref{Cliff_2c} is well-defined, note that because $\chi$ is fixed by $S_{\chi}$, $\chi \colon T \rightarrow \bC^{\times}$ can be considered as a homomorphism of $S_{\chi}$-modules with respect to conjugation on $T$ and our given action on $\bC^{\times}$. Thus $d \chi$ is just the image of $\gamma$ under of the map $ \chi \circ (-) \colon \HC^2(S_{\chi},T) \rightarrow \HC^2(S_{\chi},\bC^{\times})$.

  To see that the image of $c_{\chi}$ is $c_W$, observe that the key simplification compared to~\eqref{Cliff_1c}: the representations of $T$ are one-dimensional, so we can choose $\rho (g): g * W \to W$ to be essentially the identity maps. More precisely, for $t \in T$ and $s \in S_{\chi}$ we define
  \begin{equation} \label{Cliff_2g}
    \rho (t \breve{s}) (a) \coloneqq \chi(t)^{-1} \pi(s)(a) = \chi(t)^{-1} {}^{s}a
 =
  \begin{cases}
    \chi(t)^{-1} a , & \text{if $s$ even,}\\
    \chi(t)^{-1} \overline{a} , & \text{if $s$ odd.}\\
  \end{cases}
  \end{equation}
Then for $t_1 \breve{s_1}, t_2 \breve{s_2} \in G_{\chi}$ we can compute
\[
\rho(t_1 \breve{s_1}) \rho(t_2 \breve{s_2}) = \chi(t_1)^{-1} \pi(s_1) \circ \chi(t_2)^{-1} \pi(s_2) = \chi(t_1)^{-1} {}^{s_1}(\chi(t_2)^{-1}) \pi(s_1s_2),
\]
and
\[
\rho(t_1 \breve{s_1} t_2 \breve{s_2}) = \rho(t_1 \breve{s_1}t_2\breve{s_1}^{-1} \breve{s_1} \breve{s_2}) = \chi(t_1)^{-1} \chi( \breve{s_1}t_2\breve{s_1}^{-1})^{-1} \chi(\gamma(s_1,s_2))^{-1} \pi( s_1 s_2).
\]
Because $\chi$ is fixed by $s_2$, ${}^{s_1}(\chi(t_2)) = \chi( \breve{s_1}t_2\breve{s_1}^{-1})$, and thus
\[
d \rho(t_1 \breve{s_1}, t_2 \breve{s_2}) = \chi(\gamma(s_1,s_2))^{-1},
\]
which is exactly the pullback of $d \chi$. Independence of the choice of transversal follows from the independence of the choice of transversal of $\gamma \in \HZ^2 (S, T)$. Similarly, if $S_\chi$ is a semidirect product, the transversal is the group $S_\chi$ itself, forcing $\gamma$, and thus $d\chi$ and $c_\chi$, to be trivial.
\end{just}    
  
\begin{enumerate}[resume, label=\textnormal{(\Roman*)}]  
\item The key set of representations is ${\Irr_{\chi}^{\hspace{0.1em} \star}(S_\chi)}$, the set of irreducible finite dimensional antilinear $c_{\chi}$-projective representations of $S_\chi$. These are vector spaces $(V,\rho)$ equipped with $\pi(g)$-linear operator $\rho(g)$ for each $g\in S_\chi$ such that
\begin{equation} \label{Cliff_2d}
\chi(\gamma (h,k)) \rho (hk) = \rho(h) \rho(k), \ \mbox{ for all } \ h,k \in S_\chi \, .
\end{equation}
If $\chi(\gamma (h,k))$ is trivial, then these are just antilinear representations of $S_\chi$. In general, these can be viewed as antilinear representations of another finite group $\widetilde{S_\chi}$, a central extension of $S_\chi$ (see Section \ref{sect:projantireps}), and so we can understand ${\Irr_{\chi}^{\hspace{0.1em} \star}(S_\chi)}$ using facts about antilinear representations of finite groups.
\item Given $\chi \in T^\vee$ and $(W,\sigma)\in {\Irr_{\chi}^{\hspace{0.1em} \star}(S_\chi)}$,
the tensor product $(\bC_\chi \otimes W,\mu)$ over $\bC$ is an antilinear representation of $G_\chi$ by
\begin{equation} \label{Cliff_2e}
\mu (t \breve{s}) ( 1 \otimes w) = \chi(t) \otimes \sigma(s)(v), \mbox{ where } t\in T, s\in S_\chi, w\in W\, .
\end{equation}
\end{enumerate}

\begin{just}
Each $\mu(t \breve{s})$ is $\pi(s)$-linear, so all that needs to be checked is that $\mu(t_1 \breve{s_1}) \circ \mu(t_2 \breve{s_2}) = \mu(t_1 \breve{s_2} t_2 \breve{s_2})$, for $t_1,t_2 \in T$ and $s_1,s_2 \in S_{\chi}$. We can compute
\begin{align*}
\mu(t_1 \breve{s_1} t_2 \breve{s_2})(1 \otimes w) &= \mu(t_1 \breve{s_1} t_2 \breve{s_1}^{-1} \gamma(s_1,s_2) \breve{s_1s_2})(1 \otimes w), \\
&= \chi(t_1) \chi(\breve{s_1} t_2 \breve{s_1}^{-1}) \chi(\gamma(s_1,s_2)) \otimes \sigma(s_1s_2)(w),
\end{align*}
and 
\begin{align*}
\mu(t_1 \breve{s_1})(\mu(t_2 \breve{s_2})(1 \otimes w)) &= \mu(t_1 \breve{s_1})(\chi(t_2) \otimes \sigma(s_2)(w)), \\
&= \chi(t_1) \otimes \sigma(s_1)(\chi(t_2) \sigma(s_2)(w)), \\
&= \chi(t_2)\: {}^{s_1}(\chi(t_2)) \chi(\gamma(s_1,s_2)) \otimes \sigma(s_1 s_2)(w),\\
&= \chi(t_2)\: {}^{s_1}(\chi(t_2)) \chi(\gamma(s_1,s_2)) \otimes \sigma(s_1 s_2)(w)
\end{align*}
using that ${}^{s_1} \chi = \chi$, so ${}^{s_1}(\chi(t_2)) = \chi(\breve{s_1} t_2 \breve{s_1}^{-1})$.
\end{just}

\begin{enumerate}[resume, label=\textnormal{(\Roman*)}]  
\item
  These are all the irreducible representations.
  There is a one-to-one correspondence
\begin{equation} \label{Cliff_2f}
  \bigcup_{\chi \in S \backslash T^{\vee}}  {\Irr_{\chi}^{\hspace{0.1em} \star}(S_\chi)} \stackrel{1:1}{\longleftrightarrow}
{\Irr_{\bC}^{\hspace{0.1em} \star}(G)}, \quad
(\chi, W) \mapsto \Ind_{G_\chi}^G (\bC_\chi \otimes W , \mu)  
\end{equation}
where the union is over representatives of $S$-orbits and $(\bC_\chi \otimes W, \mu)$  is defined in~\eqref{Cliff_2e}. 
    \end{enumerate}

\begin{just}
Let $V \in \Rep_{\bC}^{\hspace{0.1em} \star}(G)$ be finite-dimensional with $V = \sum_{\chi' \in S \cdot \chi} V_{\chi'}$. By \eqref{Cliff_3a}, $\breve{r} V_{\chi} = V_{{}^{r}\chi}$ for $r \in S$, so $V = \bigoplus_{r \in S/S_{\chi}} \breve{r} V_{\chi}$, the sum being direct because eigenvectors for different characters are linearly independent. In particular, $V_{\chi}$ is stable under $G_{\chi}$, and $V \cong \Ind_{G_{\chi}}^{G} V_{\chi}$. On $W \coloneqq V_{\chi}$ the group $T$ acts by $\chi$, and setting $\sigma(s) \coloneqq \mu(\breve{s})|_{W}$ for $s \in S_{\chi}$ gives
\[
\sigma(h)\sigma(k) = \mu(\gamma(h,k)) \mu(\breve{hk}) = \chi(\gamma(h,k)) \sigma(hk),
\]
so $(W, \sigma)$ is a $c_{\chi}$-projective antilinear representation of $S_{\chi}$, and $V_{\chi} = \bC_{\chi} \otimes W$ as in \eqref{Cliff_2e}. Conversely, for any $(W, \sigma)$ the $\chi$-eigenspace of $\Ind_{G_{\chi}}^{G}(\bC_{\chi} \otimes W)$ is $\bC_{\chi} \otimes W$, as the other summands have characters ${}^{r}\chi \neq \chi$.

A $G$-stable subspace $V' \subseteq V$ is $T$-stable, hence the sum of its eigenspaces, so $V' = \bigoplus_{r} \breve{r} V'_{\chi}$ with $V'_{\chi} \subseteq W$ stable under $S_{\chi}$. Conversely an $S_{\chi}$-stable subspace $W' \subseteq W$ gives the $G$-stable subspace $\bigoplus_{r} \breve{r} W'$. Hence $V$ is irreducible if and only if $W$ is.

It follows that the map \eqref{Cliff_2f} is well-defined and that it is surjective, as every $V \in \Irr_{\bC}^{\hspace{0.1em} \star}(G)$ satisfies the hypothesis for some $\chi$ by \eqref{Cliff_2b}. It is injective, as the orbit of $\chi$ is determined by $V$ as the set of characters with $V_{\chi} \neq 0$ and $W \cong V_{\chi}$.
\end{just}
  
\begin{enumerate}[resume, label=\textnormal{(\Roman*)}]  
\item The representation $V\in {\Irr_{\bC}^{\hspace{0.1em} \star}(G)}$ is unitary if and only if the corresponding character
  $\chi\in T^\vee$ is unitary. This means that $\chi$ takes values on the unit circle $\UU(1) \leq \bC^\times$. 
    \end{enumerate} 

\begin{just}
Suppose $(V,\mu)$ is unitary. Then $\mu(t)$ is a unitary operator for each $t \in T$, so its eigenvalues lie in $\UU(1)$. This proves the only if part.

Conversely, suppose that $\chi$ is unitary, and let $V \cong \Ind_{G_{\chi}}^{G}(\bC_{\chi} \otimes W)$. As $|\chi(\gamma(h,k))| = 1$, Weyl's Unitary Trick applies to the projective representation $W$: for any positive definite hermitian form $\tau$ on $W$,
\[
\widetilde{\tau}(x,y) \coloneqq \sum_{s \in (S_{\chi})_0} \tau(\sigma(s)x, \sigma(s)y) + \sum_{s \in S_{\chi} \setminus (S_{\chi})_0} \tau(\sigma(s)y, \sigma(s)x)
\]
is preserved by $\sigma(s)$ for $s$ even and flipped for $s$ odd. As $\chi$ is unitary, the same holds for the action \eqref{Cliff_2e} of $G_{\chi}$ on $V_{\chi} = \bC_{\chi} \otimes W$. Now $V = \bigoplus_{r \in S/S_{\chi}} \breve{r} V_{\chi}$, and we define a form on $V$ by declaring the summands orthogonal and setting
\[
\langle \breve{r}u, \breve{r}v \rangle \coloneqq \begin{cases} \widetilde{\tau}(u,v), & r \text{ even}, \\ \widetilde{\tau}(v,u), & r \text{ odd}, \end{cases} \qquad u,v \in V_{\chi}.
\]
For $g \in G$ we have $g \breve{r} = \breve{r}' h$ for some $h \in G_{\chi}$, where $\pi(g) + \pi(r) = \pi(r') + \pi(h)$, and so this form is preserved by $g$ if $g$ is even and flipped if $g$ is odd.
\end{just}

\section{The Extended Stabiliser and the Dyson Type}\label{sect:extendedstabiliser}

Let $G$, $H$ and $T$ be as above, so $H$ is an index two normal subgroup of $G$ and $T$ is a finite-index normal subgroup of $G$ contained in $H$. To any character $\chi \in T^{\vee}$ and $W \in \Irr_{\chi}^{\hspace{0.1em} \star}(S_\chi)$, we saw in the last section how to obtain an irreducible antilinear representation of $G$
\[V(\chi, W) \coloneqq \Ind_{G_\chi}^G (\bC_\chi \otimes W)\, .\]

The aim of this section is to answer the question: how does the D-type of $V(\chi, W)$ depend on $(\chi,W)$? For example, if $\chi_1, \chi_2 \in T^{\vee}$ with $S_{\chi_1} = S_{\chi_2}$ and $c_{\chi_1} = c_{\chi_2} \in \HH^2(S_{\chi_1}, (\bC^{\times}, \pi))$, then for any set-map $b \colon S_{\chi_1} \rightarrow \bC^{\times}$ realising this, so $d \chi_1 = d\chi_2 db$, then there is a bijection
\[
\Omega_b \colon \Irr_{\chi_1}^{\hspace{0.1em} \star}(S_{\chi_1}) \xrightarrow{\sim} \Irr_{\chi_2}^{\hspace{0.1em} \star}(S_{\chi_2}), \qquad (W, \sigma) \mapsto (W, b\sigma), \: \text{ where }  (b\sigma)(s) \coloneqq b(s) \sigma(s).
\]
What is the relationship between the D-types of $V(\chi_1, W)$ and $V(\chi_2, \Omega_b(W))$ for $W \in \Irr_{\chi_1}^{\hspace{0.1em} \star}(S_{\chi_1})$?

It turns out that the S-types will always be the same, but  the D-types can be different. In other words, $(S_{\chi}, c_{\chi})$ does  not contain enough information to determine the D-type.

%In this section, we introduce the extended stabiliser $S_{\chi}'$ that contains enough information to remember the D-type.
%
%\subsection{The Extended Stabiliser}

The remedy is to consider {\em the extended group} $S \times C_2 = S \times \langle \epsilon \rangle$. This acts on $T^{\vee}$ by
\[
\big( (s, \epsilon^k) \cdot \chi \big)(t) \coloneqq \begin{cases} \chi(\breve{s}^{-1} t \breve{s}), & k = 0, \\ \overline{\chi(\breve{s}^{-1} t \breve{s})}, & k = 1. \end{cases}
\]

%\begin{defn}
{\em The extended stabiliser} $S'_{\chi} \leq S \times C_2$ of $\chi \in T^{\vee}$ is the stabiliser of $\chi$ for this action. It has two homomorphisms to $C_2$
\[
\pi(s, \epsilon^k) \coloneqq \pi(s) \qquad \mbox{and} \qquad \kappa(s, \epsilon^k) \coloneqq \epsilon^k.
\]
We let $S'_{\chi}$ act on $\bC^{\times}$ through $\kappa$: elements $(s,1)$ act trivially and elements $(s,c)$ act by complex conjugation. We write $\bC^{\times}_{\kappa}$ for this $S'_{\chi}$- module and define
\[
c'_{\chi} \coloneqq [d \chi']^{-1} \in \HH^2(S'_{\chi}, \bC^{\times}_{\kappa}), \qquad d\chi'\big( (s,\epsilon^j), (t, \epsilon^k) \big) \coloneqq \chi(\gamma(s,t))^{-1}.
\]
%\end{defn}

The action through $\kappa$ is the unique action of $S'_{\chi}$ on $\bC^{\times}$ by the identity and complex conjugation for which $\chi \colon T \rightarrow \bC^{\times}$ is $S'_{\chi}$-equivariant for all $\chi$. Hence, $d\chi'$ is a $2$-cocycle since it is  the image of $\gamma$ under $\chi^{-1}$. The subgroup of $S'_{\chi}$ on which $\pi = \kappa$ is $\{ (s, \epsilon^{\pi(s)}) \mid s \in S_{\chi} \}$, which we identify with $S_{\chi}$. It has index $1$ or $2$, and $d\chi'$ and $c'_{\chi}$ restrict to $d\chi$ and $c_{\chi}$ respectively on it.
We are ready to state the key fundamental fact.
%\subsection{Invariance of the D-Type}

%The statement we want to prove is the following.

To state it, consider $\chi_1, \chi_2 \in T^{\vee}$ with $S_{\chi_1}' = S_{\chi_2}'$ and $c_{\chi_1}' = c_{\chi_2}' \in \HH^2(S_{\chi_1}', \bC^{\times}_{\kappa})$.
Choose any function $b \colon S_{\chi_1}' \rightarrow \bC^{\times}$ such that $d \chi_1' = d \chi_2' db$. Since $S_{\chi_1}' = S_{\chi_2}'$, we know that $S_{\chi_1} = S_{\chi_2}$.
Consider the induced bijection
\[
\Omega_b \colon \Irr_{\chi_1}^{\hspace{0.1em} \star}(S_{\chi_1}) \xrightarrow{\sim} \Irr_{\chi_2}^{\hspace{0.1em} \star}(S_{\chi_2}), \qquad (W, \sigma) \mapsto (W, b\sigma), \: \text{ where }  (b\sigma)(s) \coloneqq b(s) \sigma(s) \, .
\]
\begin{enumerate}[label=\textnormal{(\Roman*)}]
\item  
In the notation above, the induced bijection preserves D-type. In other words, $\fD(V(\chi_1, W)) = \fD(V(\chi_2, \Omega_b(W)))$ for any $W \in \Irr_{\chi_1}^{\hspace{0.1em} \star}(S_{\chi_1})$.
\end{enumerate}
\begin{just}
  The justification takes the rest of the chapter. By Theorem~\ref{thm:dyson10}, it is sufficient to show that the associated $C_2$-graded division algebras
  \[D(\chi_1, W) \coloneqq D(V(\chi_1, W)) \quad \mbox{ and } \quad D(\chi_2, \Omega_b(W)) \coloneqq D(V(\chi_2, \Omega_b(W)))\]
  are isomorphic as $C_2$-graded algebras. We can produce an explicit isomorphism.
To do that, we describe, for a pair $(\chi, W)$, the $C_2$-graded algebra $D(\chi, W) \coloneqq D(V(\chi, W))$ purely in terms of $W$. By definition, $D(V(\chi, W)) = \End_{\bR G}(V(\chi, W))$, graded by $\bC$-linear and $\bC$-antilinear maps.

Let $\rho(t \breve{s}) \coloneqq \chi(t) \sigma(s)$ denote the action of $t \breve{s} \in G_{\chi}$ on $U \coloneqq \bC_{\chi} \otimes W \subseteq V(\chi, W)$. For $y \in S'_{\chi}$ and $\lambda \in \bC$ we write
\[{}^{y}\lambda
\coloneqq
\begin{cases}
  \lambda & \mbox{ if } \kappa(y) = 1,\\
  \overline{\lambda} & \mbox{ if } \kappa(y) = \epsilon.
\end{cases}\]
%${}^{y}\lambda$ for $\lambda$ if $\kappa(y) = 0$ and for $\overline{\lambda}$ if $\kappa(y) = 1$.
We write $\breve{y}$ for the lift of the $S$-component of $y$. If $S'_{\chi} \neq S_{\chi}$, we fix $x \in S'_{\chi} \setminus S_{\chi}$. Then $\breve{x}$ normalises $G_{\chi}$ and $\breve{x}^2 \in G_{\chi}$.
Using this element $x$, we define the even and odd components of a $C_2$-graded algebra by
%In what follows, we write
\[
D_0(U) \coloneqq \End_{\bC_{\star \hspace{-0.1em}}[G_{\chi}]}(U) = \{ \varphi_0 \colon U \rightarrow U \mid \varphi_0 \text{ is $\bC$-linear, and } \varphi_0 \, \rho(g) = \rho(g) \, \varphi_0 \text{ for all } g \in G_{\chi} \},
\]
and, if $S'_{\chi} \neq S_{\chi}$, 
\[
D_1(U) \coloneqq \{ \varphi \colon U \rightarrow U \mid \varphi(\lambda u) = {}^{x}\lambda \, \varphi(u), \text{ for all $\lambda \in \bC, u \in U$, } \varphi \, \rho(g) = \rho(\breve{x} g \breve{x}^{-1}) \, \varphi \text{ for all } g \in G_{\chi} \}.
\]
If $S'_{\chi} = S_{\chi}$, we instead set $D_1(U) \coloneqq 0$.

We need the following lemma, which we prove later.

\hrulefill
\begin{lemma}\label{lem:modelonW}
  Suppose that $\chi \in T^{\vee}$, $W \in \Irr_{\chi}^{\hspace{0.1em} \star}(S_{\chi})$, and set $V \coloneqq V(\chi, W)$. Then the map 
  \[
  D(V)_0 \rightarrow D_0(U), \qquad \Phi \mapsto \Phi|_U
  \]
  is a bijection. If $S'_{\chi} = S_{\chi}$, then $D(V)_1 = 0$. If not, then the map
  \[
  D(V)_1 \rightarrow D_1(U), \qquad \Phi \mapsto \Phi \breve{x}|_U
  \]
  is a bijection. Writing $D(U) \coloneqq D_0(U) \oplus D_1(U)$, the product on $D(V)$ transports to the product 
\[
\varphi_0 * \varphi \coloneqq \varphi_0 \circ \varphi, \qquad \varphi * \varphi_0 \coloneqq \varphi \circ \varphi_0, \qquad \varphi * \varphi' \coloneqq \varphi \circ \varphi' \circ \rho(\breve{x}^2)^{-1},
\]
on $D(U)$, where $\varphi_0 \in D_0(U)$ and $\varphi, \varphi' \in D_1(U)$.
\end{lemma}
\hrulefill

Now let $\sigma$ be the action of $S_{\chi_1}$ on $W$, so that $\Omega_b(W)$ is $W$ with the action $b\sigma$. Let us write
\[
\rho_1(t \breve{s}) = \chi_1(t) \sigma(s), \qquad \rho_2(t \breve{s}) = \chi_2(t) b(s) \sigma(s)
\]
for the two actions of $G_{\chi_1} = G_{\chi_2}$ on $W$. We apply Lemma~\ref{lem:modelonW} to both, with the same $x$ if $S'_{\chi_1} \neq S_{\chi_1}$. Let us write $D^{(1)}(W)$ and $D^{(2)}(W)$ for the resulting algebras. As $\rho_1(g)$ and $\rho_2(g)$ are proportional for all $g \in G_{\chi_1}$, we have $D^{(1)}_0(W) = D^{(2)}_0(W)$. If $S'_{\chi_1} = S_{\chi_1}$ there is nothing more to prove.
Otherwise, if $S'_{\chi_1} \neq  S_{\chi_1}$, the map
\[
D^{(1)}(W) \rightarrow D^{(2)}(W), \qquad \varphi_0 + \varphi \mapsto \varphi_0 + b(x) \varphi
\]
is an explicit isomorphism of $C_2$-graded algebras, which can be easily verified by a direct calculation.
\end{just}

Let us prove Lemma~\ref{lem:modelonW} now.

\begin{proof}
  We have that $V = \oplus_s \breve{s} U$ over the cosets of $S_{\chi}$ in $S$, and as a $T$-representation each $\breve{s} U$ is the $s \cdot \chi$-isotypical component of $V$. 
  Now for $s \in S$, $s \cdot \chi = \overline{\chi}$ if and only if $(s, \epsilon^{\pi(s) + 1}) \in S_{\chi}'$, and this happens for some $s \in S$ if and only if $S_{\chi}' \neq S_{\chi}$, as all elements of $S_{\chi}' \setminus S_{\chi}$ have this form. In particular, if $S_{\chi}' = S_{\chi}$, the $\overline{\chi}$-isotypic component of $V$ is zero, and otherwise is given by $\breve{x}^{-1} U$.

  Now as $V = \Ind_{G_{\chi}}^G(U)$, restriction is a bijection
  \[
  D(V) = \End_{\bR[G]}(V) \xrightarrow{\sim} \Hom_{\bR[G_{\chi}]}(U,V).
  \]
  For the first claim, note that this restricts to a bijection 
  \[
  D(V)_0 = \End_{\bC_{\star \hspace{-0.1em}}[G]}(V) \xrightarrow{\sim} \Hom_{\bC_{\star \hspace{-0.1em}}[G_{\chi}]}(U,V),
  \]
  and $\Hom_{\bC_{\star \hspace{-0.1em}}[G_{\chi}]}(U,V) = \Hom_{\bC_{\star \hspace{-0.1em}}[G_{\chi}]}(U,U) = D_0(U)$, because as a $T$-representation $U$ is the $\chi$-isotypic component of $V$. For the second claim, note that any $\Phi \in D(V)_1$, being antilinear, has restriction in 
  \[
  \Hom_{\bR[G_{\chi}]}(U,V_{\overline{\chi}}) \subset \Hom_{\bR[G_{\chi}]}(U,V),
  \]
  where $V_{\overline{\chi}}$ is the $\overline{\chi}$-isotypic part of $V$. Thus from above, if $S_{\chi}' = S_{\chi}$ then $\Hom_{\bR[G_{\chi}]}(U, V_{\overline{\chi}}) = 0$ and hence $D(V)_1 = 0$. If instead $S_{\chi}' \neq S_{\chi}$, then the restriction map restricts to a bijection 
  \[
  D(V)_1 \xrightarrow{\sim} \{ \psi \in \Hom_{\bR[G_{\chi}]}(U, \breve{x}^{-1} U) \mid \psi \text{ is $\bC$-antilinear}\}, \qquad \Phi \mapsto \Phi|_U.
  \]

   Now any map $\psi \colon U \rightarrow \breve{x}^{-1} U$ can be written uniquely as $\psi = \breve{x}^{-1} \varphi$ with $\varphi \colon U \rightarrow U$. As $\pi(x) \neq \kappa(x)$, $\psi$ is antilinear if and only if $\varphi(\lambda u) = {}^{x}\lambda \, \varphi(u)$ for all $\lambda \in \bC$ and $u \in U$, and $\psi$ commutes with $g \in G_{\chi}$ if and only if $\varphi \, \rho(g) = \rho(\breve{x} g \breve{x}^{-1}) \, \varphi$. Thus $\psi \mapsto \breve{x} \psi$ is a bijection from this set to $D_1(U)$, and as $\breve{x} \Phi|_U = \Phi \breve{x}|_U$, the map $\Phi \mapsto \Phi \breve{x} |_U$ is the required bijection.
  
  For the product formula, if $\Phi_0 \in D(V)_0$ and $\Phi, \Phi' \in D(V)_1$ with $\Phi_0|_U = \varphi_0 \in D_0(U)$ and $\Phi \breve{x}|_U = \varphi, \Phi' \breve{x}|_U =\varphi' \in D_1(U)$ respectively, then
  \[
  ((\Phi_0\Phi) \breve{x})|_U = \Phi_0|_U \Phi \breve{x}|_U = \varphi_0 \varphi, \qquad ((\Phi \Phi_0) \breve{x})|_U = (\Phi \breve{x}) \Phi_0|_U = \varphi \varphi_0,
  \]
  and
  \[
  (\Phi \Phi')|_U = (\Phi \breve{x} \Phi' \breve{x} \breve{x}^{-2})|_U = \varphi \varphi' \rho(\breve{x}^2)^{-1}
  \]
  as required.
\end{proof}
